\subsection{Approximation properties}

Lemma 1. --- Let X be a locally compact space, K a compact subset of X, and \((\mathrm{V}_{k})_{1\leqslant k\leqslant n}\) a finite covering of K by open sets of X. Then, there exist n continuous mappings \(f_{k}\) of X into {[}0,1{]}, such that the support of \(f_{k}\) is contained in \(V_{k}\) for \(1\leqslant k\leqslant n\) and such that \(\sum_{k=1}^{n}f_{k}(x)\leqslant1\) for all \(x\in X\) and \(\sum_{k=1}^{n}f_{k}(x)=1\) for all \(x\in K\) .

For, let \(X'\) be the compact space obtained by adjoining to X a point at infinity \(\omega\) (GT, I, §9, No.~8, Th. 4); the sets \(V_{0} = X' - K\) and \(V_{k}\) ( \(1 \leqslant k \leqslant n\) ) form an open covering of \(X'\) . Let \((f_{k})_{0 \leqslant k \leqslant n}\) be a continuous partition of unity subordinate to this covering of \(X'\) (GT, IX, §4, No.~3, Prop. 3); the functions \(f_{k}\) with index \(k \geqslant 1\) satisfy the conditions of the lemma.

Lemma 2. --- Let X be a locally compact space, K a compact subset of X, E a locally convex space, q a continuous semi-norm on E, and \(\Phi\) an equicontinuous set of mappings of X into E whose supports are contained in K. Then, for every \(\varepsilon > 0\) , there exists a continuous partition of unity \((\varphi_{j})_{0 \leqslant j \leqslant n}\) on X having the following properties:

\begin{enumerate}
\def\labelenumi{(\roman{enumi})}
\item
  \(\operatorname{Supp}(\varphi_j) \subset \mathrm{K}\) for \(1 \leqslant j \leqslant n\) .
\item
  If, for \(1 \leqslant j \leqslant n\) , \(x_j\) is any point of \(\operatorname{Supp}(\varphi_j)\) , then, for every function \(\mathbf{f} \in \Phi\) and every \(x \in \mathbf{X}\) ,
\end{enumerate}

\[
q \left(\mathbf {f} (x) - \sum_ {j = 1} ^ {n} \varphi_ {j} (x) \mathbf {f} \left(x _ {j}\right)\right) \leqslant \varepsilon .\tag{1}
\]

For every \(y\) belonging to the boundary of K, one has \(\mathbf{f}(y) = 0\) for all \(\mathbf{f} \in \Phi\) , therefore there exists an open neighborhood \(\mathrm{V}_y\) of \(y\) in X such that, for every \(z \in \mathrm{V}_y\) and every \(\mathbf{f} \in \Phi\) , one has \(q(\mathbf{f}(z)) \leqslant \varepsilon / 2\) . Let \(\mathrm{K}'\) be the set of points of K not belonging to any of the \(\mathrm{V}_y\) as \(y\) runs over the boundary of K; \(\mathrm{K}'\) is compact and is contained in the interior of K. The set \(\Phi\) is uniformly equicontinuous in K; therefore, there exists a finite open covering \((\mathrm{U}_j)_{1 \leqslant j \leqslant n}\) of \(\mathrm{K}'\) consisting of open sets in X contained in K, such that for every pair of points \(x, y\) of a same \(\mathrm{U}_j\) , one has \(q(\mathbf{f}(x) - \mathbf{f}(y)) \leqslant \varepsilon / 2\) for every \(\mathbf{f} \in \Phi\) . By Lemma 1, there exist \(n\) continuous mappings \(\varphi_j\) of X into {[}0,1{]} ( \(1 \leqslant j \leqslant n\) ) such that \(\operatorname{Supp}(\varphi_j) \subset \mathrm{U}_j\) and such that \(\sum_{j=1}^{n} \varphi_j(x) \leqslant 1\) on X and \(\sum_{j=1}^{n} \varphi_j(x) = 1\) on \(\mathrm{K}'\) . For \(x_j \in \operatorname{Supp}(\varphi_j)\) ( \(1 \leqslant j \leqslant n\) ) and \(\mathbf{f} \in \Phi\) , we therefore have, for all \(x \in \mathrm{U}_j\) ,

\[
q \big (\mathbf {f} (x) \varphi_ {j} (x) - \mathbf {f} (x _ {j}) \varphi_ {j} (x) \big) = \varphi_ {j} (x) q \big (\mathbf {f} (x) - \mathbf {f} (x _ {j}) \big) \leqslant \frac {\varepsilon}{2} \varphi_ {j} (x),
\]

and this relation remains true if \(x \notin U_{j}\) since then \(\varphi_{j}(x) = 0\) . By addition we infer that, for every \(x \in X\) ,

\[
q \left(\mathbf {f} (x) \left(1 - \varphi_ {0} (x)\right) - \sum_ {j = 1} ^ {n} \varphi_ {j} (x) \mathbf {f} \left(x _ {j}\right)\right) \leqslant \frac {\varepsilon}{2} \left(1 - \varphi_ {0} (x)\right),\tag{2}
\]

where \(\varphi_0 = 1 - \sum_{j=1}^{n} \varphi_j\) ; whence (1) for \(x \in \mathbf{K}'\) , since then \(\varphi_0(x) = 0\) ; (1) also holds for \(x \notin \mathbf{K}\) , the first member then being zero. Finally, for \(x \in \mathbf{K} - \mathbf{K}'\) we have \(q(\mathbf{f}(x)\varphi_0(x)) \leqslant \varepsilon/2\) by the definition of \(\mathbf{K}'\) , therefore this relation and (2) again imply (1) in this case.

Let X be a locally compact space; for every Banach space E (real or complex) we denote by \(\mathcal{C}^{b}(X;E)\) the vector space of continuous and bounded mappings of X into E; we know that the topology of uniform convergence in X is compatible with the vector space structure (real, resp. complex) of \(\mathcal{C}^{b}(X;E)\) , and it is defined by the norm

\[
\| \mathbf {f} \| = \sup _ {x \in X} \| \mathbf {f} (x) \|.\tag{3}
\]

Moreover, the normed space thus defined is a Banach space (GT, X, §3, No.~2, and No.~1, Cor. 2 of Prop. 2); the topology defined by this norm on \(\mathcal{K}(\mathrm{X};\mathrm{E})\) (in other words, the topology of uniform convergence in X) is coarser than the direct limit topology on \(\mathcal{K}(\mathrm{X};\mathrm{E})\) defined in No.~1.

PROPOSITION 3. --- Let X be a locally compact space, X' the compact space obtained by adjoining to X a point at infinity ω (GT, I, §9, No.~8,

Th. 4), and E a Banach space. The closure of \(\mathcal{K}(\mathrm{X};\mathrm{E})\) in the normed space \(\mathcal{C}^b (\mathrm{X};\mathrm{E})\) is the vector space of continuous functions on \(\mathbf{X}\) , with values in E and tending to 0 at the point \(\omega\) .

Let \(\mathbf{f} \in \mathcal{C}^b(\mathrm{X};\mathrm{E})\) be a function in the closure of \(\mathcal{K}(\mathrm{X};\mathrm{E})\) ; for every \(\varepsilon > 0\) , there exists a function \(\mathbf{g} \in \mathcal{K}(\mathrm{X};\mathrm{E})\) such that \(\| \mathbf{f}(x) - \mathbf{g}(x) \| \leqslant \varepsilon\) for all \(x \in \mathrm{X}\) ; if \(\mathrm{K}\) is the support of \(\mathbf{g}\) , it follows that \(\| \mathbf{f}(x) \| \leqslant \varepsilon\) for all \(x \in \complement \mathrm{K}\) , thus \(\mathbf{f}(x)\) tends to 0 as \(x\) tends to \(\omega\) . Conversely, if \(\mathbf{f}\) has this property then, for every \(\varepsilon > 0\) , there exists a compact set \(\mathrm{K} \subset \mathrm{X}\) such that \(\| \mathbf{f}(x) \| \leqslant \varepsilon\) for all \(x \in \complement \mathrm{K}\) . By Lemma 1 there exists a continuous mapping \(h\) of X into {[}0,1{]}, with compact support, equal to 1 on \(\mathrm{K}\) ; then \(\| \mathbf{f}(x) h(x) \| \leqslant \varepsilon\) on \(\complement \mathrm{K}\) and \(\mathbf{f}(x) = \mathbf{f}(x) h(x)\) on \(\mathrm{K}\) ; since \(\mathbf{f}h\) has compact support and \(\| \mathbf{f}(x) - \mathbf{f}(x) h(x) \| \leqslant 2\varepsilon\) for all \(x \in \mathrm{X}\) , the proposition is proved.

We shall denote by \(\mathcal{C}^{0}(\mathrm{X};\mathrm{E})\) the subspace of \(\mathcal{C}^{b}(\mathrm{X};\mathrm{E})\) formed by the functions tending to zero at the point at infinity \(\omega\) ; it is thus the completion of the normed space \(\mathcal{K}(\mathrm{X};\mathrm{E})\) .

PROPOSITION 4. --- Let X be a locally compact space, E a locally convex space; then, the space \(\mathcal{K}(\mathrm{X};\mathrm{E})\) is dense in \(\mathcal{C}(\mathrm{X};\mathrm{E})\) for the topology of compact convergence.

For every compact set \(K \subset X\) , there exists a function \(h \in \mathcal{K}(X; \mathbf{R})\) equal to 1 on K, by Lemma 1; for every function \(\mathbf{f} \in \mathcal{C}(X; E)\) the function hf, which belongs to \(\mathcal{K}(X; E)\) , is equal to f on K, whence our assertion.

PROPOSITION 5. --- Let X be a locally compact space, E a real (resp. complex) locally convex space. For every compact subset K of X, the vector space \(\mathcal{K}(X,K;\mathbf{R})\otimes_{\mathbf{R}}E\) (resp. \(\mathcal{K}(X,K;\mathbf{C})\otimes_{\mathbf{C}}E\) ) (identified with a set of mappings of X into E, cf.~A, II, §7, No.~7, Cor. of Prop. 15) is dense in \(\mathcal{K}(X,K;E)\) ; the vector space \(\mathcal{K}(X;\mathbf{R})\otimes_{\mathbf{R}}E\) (resp. \(\mathcal{K}(X;\mathbf{C})\otimes_{\mathbf{C}}E\) ) is dense in \(\mathcal{K}(X;E)\) .

As the second assertion is an obvious consequence of the first, it suffices to prove the latter. We apply Lemma 2 with \(\Phi\) reduced to a single element f of \(\mathcal{K}(X,K;E)\) ; then, for every \(x\in X\) ,

\[
q \left(f (x) - \sum_ {j = 1} ^ {n} \varphi_ {j} (x) f \left(x _ {j}\right)\right) \leqslant \varepsilon ,
\]

where the \(\varphi_{j}\) belong to \(\mathcal{K}(\mathrm{X},\mathrm{K};\mathbf{R})\) ; since the mapping \(x\mapsto \sum_{j = 1}^{n}\varphi_{j}(x)f(x_{j})\) may be canonically identified with the element \(\sum_{j = 1}^{n}\varphi_{j}\otimes f(x_{j})\) , this proves the proposition, by the definition of the topology of \(\mathcal{K}(\mathrm{X},\mathrm{K};\mathrm{E})\) .

